\begin{abstract}
  \noindent\textbf{Abstract}\par\vspace{0.5em}
  With the desire to react quickly to events or complete time-sensitive tasks,
  IoT devices often need real-time data interpretation. The design and modeling
  of a fast and efficient triple-cascode operational transconductance amplifier
  (TC-OTA) that operate rapidly and effectively deliver rapid signal
  conditioning and amplification, enabling the processing of real-time sensor
  data is presented in this study. Carbon nanotube field-effect transistors
  (CNTFETs) and metal oxide semiconductor field-effect transistor (MOSFET) are
  used to implement the suggested TC-OTAs at 32 nm technology node. A
  traditional CMOS-based TC-OTA is comprehensively compared with three distinct
  CNT-based TC-OTA architectures. A thorough analysis of key performance
  indicators reveals that the CNT-based designs have undergone significant
  improvements, with the pure CNT-TC-OTA exhibiting the most substantial
  enhancement. In particular, the pure CNT-TC-OTA has a 197\% higher DC gain
  than the CMOS-TC-OTA, a 4.61 times improvement in the common-mode rejection
  ratio (CMRR), and a 38.08\% decrease in power consumption. Additionally, the
  suggested pure CNT-TC-OTA improves figures of merit (FOM1 and FOM2) by 4.53
  times and 1.21 times, respectively. The impacts of three important CNT
  parameters—nanotube diameter (DCNT), pitch (S), number of CNTs (N), oxide
  thickness (TOX) and dielectric constant (KOX)—on the OTA performance are
  methodically and thoroughly examined at a load capacitance of 0.01 pF and a
  supply voltage of 0.9 V. The findings suggest that selecting the optimal CNT
  diameter, pitch, and count can further enhance circuit performance.
  Furthermore, when compared to the traditional CMOS-based TC-OTA, the pure
  CNT-TC-OTA exhibits noticeably better stability characteristics. Furthermore,
  the proposed CNT-based triple cascode OTA has been compared with an identical
  OTA circuit at a 45 nm technology node to observe the benefits of a 32 nm
  technology node. Significant improvements in gain (54.82\%), bandwidth
  (24.81\%), slew rate (162\%), and CMRR (46.11\%) have been reported, along
  with a 35.33\% decrease in average power, when the channel length in the
  CNTFET-based OTA is shortened.
\end{abstract}